\subsection{CPython Dynamic Scaling Mechanics}

A \texttt{PyListObject} tracks two quantities: \texttt{ob\_size} (logical length, exposed as \texttt{len()}) and allocated slot capacity (implementation field, typically larger). Appends that fit in spare capacity write one new pointer without resizing; capacity exhaustion triggers reallocation.

\subsubsection{Dynamic Over-Allocation Invariants: How CPython Pre-Allocates Extra Slots During List Resizing to Support Amortized O(1) Appends}

When \texttt{list\_resize} grows the pointer array, CPython does not allocate exactly \texttt{newsize} slots. The internal growth rule (CPython C source) is:

\begin{verbatim}
new_allocated = newsize + (newsize >> 3) + (newsize < 9 ? 3 : 6)
\end{verbatim}

Extra headroom (\texttt{newsize >> 3} plus a small constant) amortizes append cost: most appends become pointer writes into pre-reserved slots; occasional appends hit a threshold, allocate a larger \texttt{PyObject **} block, copy all existing pointers, and free the old array---a latency spike visible as a step in \texttt{sys.getsizeof()}. Growth therefore alternates between cheap slot assignment and expensive bulk pointer copying.

\subsubsection{Memory Shifting Costs: The O(N) Penalty of Arbitrary Index Insertions and Deletions (\texttt{.insert()}, \texttt{.pop()})}

Tail \texttt{append}/\texttt{pop} avoid shifting existing slots. \texttt{insert(0, x)} or \texttt{pop(0)} moves every surviving pointer one position---$O(n)$ pointer writes per operation. Repeated front insertion on a large list is structurally quadratic; \texttt{collections.deque} uses a block-linked layout for $O(1)$ ends at the cost of worse middle indexing.
